\section{Discussion}\label{sec:discussion}

\paragraph{Relatedness and tokenizer fit} Relatedness is the most consistent predictor of transfer in this design, and about a third of its effect is shared with the tokenizer-fertility gap. Fertility helps to explain pairs that relatedness alone does not, such as Swahili transferring well into the Sotho-Tswana languages. However, the two are partly collinear here, therefore we treat fertility as a candidate mechanism behind part of the relatedness effect and not as a replacement for it.

\paragraph{What ``domain match'' measures} The corpus effect is concentrated in isiZulu and in news targets: a same-corpus sister language beats the same language from the other corpus by 32.0~F1 on isiZulu news, while on isiXhosa and Setswana government text the same-language news source wins. NCHLT and MasakhaNER differ in text domain, but they also differ in annotation guidelines and annotation teams. Thus, our design cannot tell whether the corpus effect comes from the text or from how each corpus defines an entity, for example whether titles and role words belong inside person spans. This distinction matters in practice because labelling conventions can be changed while domain cannot, and Section~\ref{sec:future} proposes experiments that separate the two.

\paragraph{Exposure and the Setswana reversal} Documented pretraining exposure showed no detectable association; however, with three seen and three unseen targets that line up with the language clusters, the design has little power to find one. For Setswana, Swahili outperforms both Sotho-Tswana relatives, and an equal number of entity spans does not change this. Since the relatives come only from government text and Swahili comes from news, the reversal is again an effect of source language and corpus together.

\paragraph{Relation to prior work} The advantage of news sources over government sources agrees with MphayaNER, where related-language news sources beat Tshivenda government data~\cite{mbuvha2023}, and with the low government-domain zero-shot scores in NGLUEni~\cite{meyer2024}. Furthermore, our Setswana ranking extends the MasakhaNER~2.0 LangRank analysis~\cite{lin2019,adelani2022} with the two Sotho-Tswana relatives that it could not include, and as annotated, both rank below Swahili.

\paragraph{Practical implications} For news NER in these languages, a news-annotated related source is a better starting point than same-language government data, and Swahili is a strong Setswana source. One hundred target sentences help the most when the only available source comes from the other corpus, and 500 sentences bring every source to within 1--5~F1 of the in-language model.
